\problemname{Löjtnant Slotto}
Löjtnant Slotto befäster just nu $N$ stycken slott.
Varje slott har rikedomar värda $l_i$ guldmynt, men vaktas av $g_i$ soldater.
Du vill nu anfalla Slottos slott.
Till ditt förfogande har du $T$ trupper.
Varje trupp kan skickas för att anfalla exakt ett slott.
Om du skickar $g_i$ eller mer soldater till slott $i$, så kan du plundra alla dess rikedomar.
Annars kommer anfallet att misslyckas.

Om du fördelar dina trupper för anfall optimalt, hur många guldmynt värt av rikedomar kan du få?

\section*{Indata}
Den första raden av indata innehåller heltalen $N$ och $T$ ($1 \leq N \leq 2 \cdot 10^5$, $1 \leq T \leq 10^6$),
antalet slott det finns och antalet trupper du har.

De följande $N$ raderna beskriver vardera ett slott.
Varje av dessa rader innehåller heltalen $l_i$ och $g_i$ ($1 \leq l_i \leq 5$, $1 \leq g_i \leq 10^6$),
värdet av rikedomarna i slott $i$ och antalet vaktar som vaktar slott $i$.

\section*{Utdata}
Skriv ut ett heltal: största mängden rikedomar du kan erövra om du fördelar dina trupper optimalt.

\section*{Poängsättning}
Din lösning kommer att testas på en mängd testfallsgrupper.
För att få poäng för en grupp så måste du klara alla testfall i gruppen.

\noindent
\begin{tabular}{| l | l | p{12cm} |}
  \hline
  \textbf{Grupp} & \textbf{Poäng} & \textbf{Gränser} \\ \hline
  $1$    & $7$   & $l_i=1$ \\ \hline
  $2$    & $12$  & $N, T \leq 500$ \\ \hline
  $3$    & $19$  & $N \leq 2000$ \\ \hline
  $4$    & $22$  & $l_i \leq 2$ \\ \hline
  $5$    & $40$  & Inga ytterligare begränsningar. \\ \hline
\end{tabular}

